\chapter{Diario de desarrollo}
\label{anx:diario}

Registro cronológico de las sesiones de trabajo: tarea realizada, resultado y decisiones tomadas. Las decisiones técnicas con alternativas se documentan además como \gls{adr} en el directorio \texttt{docs/adr} del repositorio.

\begin{longtable}{@{}>{\raggedright\arraybackslash}p{2.1cm}>{\raggedright\arraybackslash}p{4.2cm}>{\raggedright\arraybackslash}p{7.7cm}@{}}
\caption{Diario de desarrollo.}\label{tab:diario}\\
\toprule
\textbf{Fecha} & \textbf{Tarea} & \textbf{Resultado y decisiones} \\
\midrule
\endfirsthead
\toprule
\textbf{Fecha} & \textbf{Tarea} & \textbf{Resultado y decisiones} \\
\midrule
\endhead
\bottomrule
\endfoot
10/11/2026 & Aprobación de la propuesta & El tutor da el visto bueno a la propuesta. Se acuerda ajustar el alcance y las funcionalidades según el tiempo disponible y los avances. \\
10/11/2026 & Definición del proyecto & Se fija el \textit{stack} (Python, Dash, SQLite, pip/venv), las convenciones del repositorio y las \textit{skills} de los asistentes de \gls{ia}. ADR 0001--0003. \\
10/11/2026 & Plantilla de la memoria & Se adopta la plantilla de TFG de la EPSC (v3.13) dentro del repositorio. El DNI del Anexo V se mantiene fuera del repositorio público. \\
10/11/2026 & Estructura de la memoria & Capítulos organizados según las partes de la propuesta; se eliminan las secciones de la plantilla que no aplican a un TFG de software. Al verificar la bibliografía se corrigen dos referencias: Davenport \textit{et al.} (2010, no 2018) y Géron (2022, no 2023). \\
12/11/2026 & Organización del proyecto & Se separa la planificación del TFG en un espacio propio de Notion y se fija una única fuente de verdad por tipo de información (decisiones en ADR, convenciones en AGENTS.md, planificación en Notion). Se reconfirma el \textit{stack} del ADR 0001 frente a referencias antiguas a FastAPI y PostgreSQL. \\
12/11/2026 & Fuente de datos e ingesta & Los únicos datos reales están en adidas Running, cuya \gls{api} está cerrada desde 2025. Se decide leer su exportación \gls{json} original e implementar la ingesta con el patrón adaptador hacia un esquema canónico, para que cambiar de fuente no afecte al resto del sistema. ADR 0004. Pendiente: comprobar si la exportación incluye \gls{fc}. \\
12/11/2026 & Revisión de la fuente de datos & Se dispone de la exportación de Apple Salud, que reúne todas las carreras (128, con 17 duplicadas entre aplicaciones). Se cambia la fuente del primer adaptador a esta exportación y se fija una fuente única, los entrenamientos del Apple Watch (71 carreras, 70 con \gls{fc} media), para no mezclar sensores. ADR 0004 revisado antes de su fusión. \\
\end{longtable}
